\chapter*{Reading guide}
\addcontentsline{toc}{chapter}{Reading guide}

The construction has two limiting operations. First the lattice spacing
is removed while the physical volume is held fixed. Then the volume
increases along a subsequence. The first operation requires a stability
argument, not just compactness. The second uses compactness only after
bounds uniform in volume have been established. Keeping these operations
separate also fixes the meaning of the local observables throughout the
book.

A lattice expression with the correct classical limit is not yet a
normalized quantum field. Integration can mix it with other expressions
of the same or lower degree, including the constant. We therefore begin
with finite-dimensional spaces of polynomial labels, construct their
local lattice representatives, and later determine their normalization
by transporting all coefficients together. Fields are finally defined
by equality of all separated correlation responses. This prevents a
formal equation of motion from being identified prematurely with a
vanishing interacting field.

The main analytic work is the block step. Its quadratic part determines
a Gaussian fluctuation field and a stationary background. The nonlinear
integral is separated into regions of small curvature, where Taylor and
Gaussian estimates apply, and complementary regions, where a positive
action cost is retained. Local terms are reorganized as polymer
activities: functions associated with connected supports, controlled by
weighted norms. The support weights must accommodate the full dependence
of an activity, including the cutoffs and local gauge choices used to
define it. This distinction matters when several insertions interact.

Once the block estimates are available, one scale schedule must satisfy
all fixed finite requirements. It is chosen slowly enough that each
stage's constants can be controlled before the next set of requirements
is introduced. A common trajectory then determines the couplings,
activities and observable normalizations for the entire Wilson family.
Observable insertions are tracked by commuting auxiliary variables whose
squares are zero. Coefficients involving several variables retain the
terms generated when insertions merge under blocking.

The continuum part uses this information in a different order from a
formal perturbative calculation. Finite families of free-field correlation functionals separate the
polynomial labels and control the dimension of the interacting field spaces. Coordinate
reflection positivity constructs Hilbert spaces and bounded slab
operators before continuous rotations are known. A conserved tensor,
local flux estimates and positivity then establish rotational Ward
identities and fix the stress normalization. Only then are phases
separated and relativistic fields reconstructed.

\begin{center}
\begin{tabular}{@{}p{.17\textwidth}p{.76\textwidth}@{}}
\toprule
Chapters & Role in the argument\\
\midrule
1--3 & Lattice conventions, local labels and exact microscopic equations;
free Maxwell correlations that separate polynomial labels.\\[5pt]
4--9 & Exact blocking, Gaussian kernels, large-field estimates,
local matching, polymer expansions and the nonlinear step.\\[5pt]
10--12 & One common trajectory, observable insertions and the full
ultraviolet limit at fixed physical volume.\\[5pt]
13--16 & Infinite-volume correlations, field quotients, Ward identities,
phases and relativistic reconstruction.\\[5pt]
17 & The quantitative trace coefficient and its all-radius consequence.\\[5pt]
18--19 & Finite operator products, specified contact extensions,
and assembly with the exact scope of the conclusions.\\
\bottomrule
\end{tabular}
\end{center}

The final two analytic refinements have different purposes. The trace
law needs a rate beyond convergence to the traceless free tensor; an
error that tends to zero may still be too large after multiplication by
the running coupling. An operator product expansion needs a signed
replacement estimate for a whole cluster, uniform under repeated copies.
Neither follows from the existence of separated correlations alone.
A conservative finite-order extension across collision diagonals,
by contrast, follows from polynomial separated bounds and a specified
Taylor subtraction.

The final assembly chapter collects the proved dependencies and
distinguishes the conclusions from further questions not addressed here. The
accompanying chapter audits provide source locations and review records;
they are working documents rather than prerequisites for reading the
mathematical exposition.
